%%% LaTeX-Vorlage Version 2.1 %%%

% TODO: individuelle Einstellungen (Name, Titel etc.)
% -> bitte in Konfigurationsdatei anpassen
%%%%
%
% Zentrale Konfigurationsdatei
%
% In dieser Datei sind eine Reihe verpflichtender Einstellungen 
% (Nr. 1 bis 6) vorzunehmen.
%
% Die Einstellungen unter Nr. 7 bis 11 können im Regelfall unverändert
% belassen werden. Ausnahmen sind:
%  - Ihre Arbeit ist in englischer Sprache verfasst (Nr. 7)
%  - der Titel Ihrer Arbeit ist sehr lang, so dass er nicht auf das
%    Deckblatt passt oder anders umgebrochen werden soll (Nr. 8 und 9)
%  - es soll ein besonderes Abgabedatum angegeben werden (Nr. 10)
%  - Sie benötigen einen Vertraulichkeitsvermerk (Nr. 11)
%
%%%%


% TODO 1. Typ der Arbeit (für Titelseite und Metadaten)
% Zutreffendes auswählen:

\newcommand{\typMeinerArbeit}{Projektarbeit T2000} 
%\newcommand{\typMeinerArbeit}{PA1} 
%\newcommand{\typMeinerArbeit}{PA2} 
%\newcommand{\typMeinerArbeit}{Seminararbeit} 
%\newcommand{\typMeinerArbeit}{BA} 

% TODO 2. Vorname, Name der Autorin/des Autors (für Deckblatt und Metadaten)
\newcommand{\meinName}{Tudor Lupsa}

% TODO 3. Kurs eintragen
\newcommand{\meinKurs}{xxx}

% TODO 4. Titel der Arbeit (für Deckblatt, ehrenwörtliche Erklärung und Metadaten, ohne Umbrüche angeben)
\newcommand{\themaMeinerArbeit}{Mein Titel}

% TODO 5. Angaben zum Unternehmen (wird bei Seminararbeiten nicht angezeigt)
\newcommand{\UNName}{Firma XY}
\newcommand{\UNBetreuer}{Titel, Vorname und Nachname}
\newcommand{\UNBetreuerFunktion}{Funktion des Betreuers/der Betreuerin}

% TODO 6. Angaben zur wissenschaftlichen Betreuung 
\newcommand{\DHBWBetreuer}{Titel, Vorname und Nachname}


% OPTIONALE Einstellungen

% 7. Arbeit in Englisch
% (nur ändern, falls Ihre Arbeit in englischer Sprache geschrieben ist)
\newcommand{\meineSprache}{DE}	% Standard-Einstellung
%\newcommand{\meineSprache}{EN}	% für Arbeiten in englischer Sprache

% 8. Schriftgröße des Titels auf Deckblatt
% (nur ändern, falls Sie einen sehr langen Titel haben)
% Zutreffendes auswählen:
\newcommand{\schriftgroesseTitel}{\LARGE}   % Standard-Einstellung
%\newcommand{\schriftgroesseTitel}{\Large}  % bei sehr langen Titeln

% 9. Titel mit Umbrüchen für Deckblatt
% (nur ändern, falls Sie den Zeilenumbruch selbst beeinflussen möchten)
\newcommand{\titelAufDeckblatt}{\themaMeinerArbeit}		% Standard-Einstellung
%\newcommand{\titelAufDeckblatt}{Herausforderungen der Digitalisierung im globalen Wettbewerb von Industrieunternehmen \\ -- eine vergleichende Untersuchung unter Berücksichtigung aktueller und weniger aktueller Forschungsmethoden \\ am Beispiel der Firma Melanie Müller und Söhne AG} % explizite Angabe

% 10. Abgabedatum anpassen
% Zutreffendes auswählen:
\newcommand{\abgabeDatum}{\today}  		% Standard-Einstellung
%\newcommand{\abgabeDatum}{TT.MM.JJJJ}  % falls nicht aktuelles Datum

% 11. Vertraulichkeitsvermerk
% (nur ändern, falls Ihre Arbeit einen Vertraulichkeitsvermerk tragen soll)
\newcommand{\hatVermerk}{nein}  	% Standard-Einstellung
% \newcommand{\hatVermerk}{ja}  	% falls Vertraulichkeitsvermerk



% Grundlegende Dokumenteneigenschaften gemäß DHBW-Vorgaben
\documentclass[a4paper,fontsize=11pt,oneside,parskip=half,headings=normal,listof=nochaptergap]{scrreprt} 
% \usepackage{showframe} % nur für Kontrolle der Ränder 

%%% Präambel einbinden (mit Festlegungen gemäß DHBW-Vorgaben) %%%
%%% Präambel %%%
% Sprachumschaltung DE/EN
\usepackage{ifthen}
\newcommand{\DEoEN}[2]{\ifthenelse{\equal{\meineSprache}{DE}}{#1}{#2}}

% Zeichencodierung/Fonts
\usepackage[utf8]{inputenc}
\usepackage[T1]{fontenc}

% Farben + Code-Listings
\usepackage{xcolor}      % statt 'color' (mehr Features)
\usepackage{listings}
% Grund-Setup für Listings (Umlaute/EUR etc.)
\lstset{numbers=left, numberstyle=\tiny, numbersep=5pt, texcl=true}
\lstset{literate=
{Ö}{{\"O}}1 {Ä}{{\"A}}1 {Ü}{{\"U}}1 {ß}{{\ss}}2
{ü}{{\"u}}1 {ä}{{\"a}}1 {ö}{{\"o}}1 {€}{{\euro}}1
}
% Lesbares Listing-Layout + automatischer Zeilenumbruch
\lstdefinestyle{code}{
  numbers=left, numberstyle=\tiny, numbersep=5pt,
  breaklines=true, breakatwhitespace=true,
  columns=fullflexible, keepspaces=true, tabsize=2,
  basicstyle=\ttfamily\small,
  postbreak=\mbox{\textcolor{gray}{$\hookrightarrow$}}\space
}
\lstset{style=code}

% Seitenränder
\usepackage[
  left=2.5cm, right=2.5cm,
  top=2.5cm, bottom=2.5cm,
  foot=12mm, includefoot
]{geometry}

% Sprache + Anführungszeichen
\DEoEN{
  \usepackage[ngerman]{babel}
  \usepackage[babel,german=quotes]{csquotes}
}{
  \usepackage[english]{babel}
  \usepackage[babel,english=british]{csquotes}
}

% Listen, Grafiken, Zeilenabstand
\usepackage{enumerate}
\usepackage{graphicx}
\graphicspath{{img/}}
\usepackage[onehalfspacing]{setspace}

% Testtext, Akronyme
\usepackage{blindtext}
% \usepackage{color}  % durch xcolor ersetzt
\usepackage[nohyperlinks]{acronym}

% Literatur (biblatex + DHBW-Config)
\usepackage[
  backend=biber,
  bibstyle=_dhbw_authoryear,maxbibnames=99,
  citestyle=authoryear, dashed=false,
  uniquename=true, useprefix=true,
  bibencoding=utf8
]{biblatex}
%kein Punkt am Ende bei \footcite
%http://www.golatex.de/footcite-ohne-punkt-am-schluss-t4865.html
\renewcommand{\bibfootnotewrapper}[1]{\bibsentence#1}

% Bibliographie: Vornamen ausgeschrieben
\DeclareNameAlias{author}{family-given}
\DeclareNameAlias{editor}{family-given}

%Reihenfolge der Autorennamen
%   
% http://golatex.de/viewtopic,p,80448.html#80448
% Argumente: siehe http://texwelt.de/blog/modifizieren-eines-biblatex-stils/
\DeclareNameFormat{sortname}{% Bibliographie
  \ifnum\value{uniquename}=0 % Normalfall
    \ifuseprefix%
      {%
         \usebibmacro{name:family-given}
           {\namepartfamily}
           {\namepartgiveni}
           {\namepartprefix}
           {\namepartsuffixi}%
       }
      {%
         \usebibmacro{name:family-given}
           {\namepartfamily}
           {\namepartgiveni}
           {\namepartprefixi}
           {\namepartsuffixi}%
       }%
  \fi
  \ifnum\value{uniquename}=1% falls nicht eindeutig, abgek. Vorname 
      {%
         \usebibmacro{name:family-given}
           {\namepartfamily}
           {\namepartgiveni}
           {\namepartprefix}
           {\namepartsuffix}%
       }%
  \fi
  \ifnum\value{uniquename}=2% falls nicht eindeutig, ganzer Vorname 
      {%
         \usebibmacro{name:family-given}
           {\namepartfamily}
           {\namepartgiven}
           {\namepartprefix}
           {\namepartsuffix}%
       }%
  \fi   
  \usebibmacro{name:andothers}}

\DeclareNameFormat{labelname}{% für Zitate
  \ifnum\value{uniquename}=0 % Normalfall
    \ifuseprefix%
      {%
         \usebibmacro{name:family-given}
           {\namepartfamily}
           {\empty}
           {\namepartprefix}
           {\namepartsuffixi}%
       }
      {%
         \usebibmacro{name:family-given}
           {\namepartfamily}
           {\empty}
           {\namepartprefixi}
           {\namepartsuffixi}%
       }%
  \fi
  \ifnum\value{uniquename}=1% falls nicht eindeutig, abgek. Vorname 
      {%
         \usebibmacro{name:family-given}
           {\namepartfamily}
           {\namepartgiveni}
           {\namepartprefix}
           {\namepartsuffix}%
       }%
  \fi
  \ifnum\value{uniquename}=2% falls nicht eindeutig, ganzer Vorname 
      {%
         \usebibmacro{name:family-given}
           {\namepartfamily}
           {\namepartgiven}
           {\namepartprefix}
           {\namepartsuffix}%
       }%
  \fi   
  \usebibmacro{name:andothers}}
      
  
\DeclareFieldFormat{extrayear}{% = the 'a' in 'Jones 1995a'
  \iffieldnums{labelyear}
    {\mknumalph{#1}}
    {\mknumalph{#1}}}        

% Namen getrennt durch Komma (Zitate)
\DeclareDelimFormat*[footcite,cite,textcite,parencite]{multinamedelim}{\addcomma\space}
% bzw. Semikolon (Literaturverzeichnis)
\DeclareDelimFormat[bib,biblist]{multinamedelim}{\addsemicolon\space}
% keine besondere Behandlung beim letzten Autor
\DeclareDelimAlias{finalnamedelim}{multinamedelim}
%\DeclareDelimAlias{multilistdelim}{multinamedelim}

\renewcommand{\nameyeardelim}{~}

% Literaturverzeichnis: Doppelpunkt zwischen Name (Jahr): Rest 
% http://de.comp.text.tex.narkive.com/Tn1HUIXB/biblatex-authoryear-und-doppelpunkt
\renewcommand{\labelnamepunct}{\addcolon\addspace}

% damit die Darstellung für Vollzitate von Primärquellen in 
% Fußnoten später auf "nicht fett" geändert werden kann 
% (nur für Zitate von Sekundärliteratur relevant)
\newcommand{\textfett}[1]{\textbf{#1}}

% für Zitate von Sekundärliteratur:
\newcommand{\footcitePrimaerSekundaer}[4]{%
  \renewcommand{\textfett}[1]{##1}%
  \footnote{\fullcite[#2]{#1}, \DEoEN{zitiert nach}{as cited in} \cite[#4]{#3}}%  
  \renewcommand{\textfett}[1]{\textbf{##1}}%
}

% Im Literaturverzeichnis: Autor (Jahr) fett
\renewbibmacro*{author}{%
  \ifboolexpr{%
    test \ifuseauthor%
    and
    not test {\ifnameundef{author}}
  }
    {\usebibmacro{bbx:dashcheck}
       {\bibnamedash}
       {\usebibmacro{bbx:savehash}%
        \textfett{\printnames{author}}%
        \iffieldundef{authortype}
          {\setunit{\addspace}}
          {\setunit{\addcomma\space}}}%
     \iffieldundef{authortype}
       {}
       {\usebibmacro{authorstrg}%
        \setunit{\addspace}}}%
    {\global\undef\bbx@lasthash
     \usebibmacro{labeltitle}%
     \setunit*{\addspace}}%
  \textfett{\usebibmacro{date+extrayear}}}

% Sonderfall: Quelle ohne Autor, aber mit Herausgeber
% Name des Herausgebers wird fett gedruckt
\renewbibmacro*{bbx:editor}[1]{%
  \ifboolexpr{%
    test \ifuseeditor%
    and
    not test {\ifnameundef{editor}}
  }
    {\usebibmacro{bbx:dashcheck}
       {\bibnamedash}
       {\textfett{\printnames{editor}}%
        \setunit{\addcomma\space}%
        \usebibmacro{bbx:savehash}}%
     \usebibmacro{#1}%
     \clearname{editor}%
     \setunit{\addspace}}%
    {\global\undef\bbx@lasthash
     \usebibmacro{labeltitle}%
     \setunit*{\addspace}}%
  \textfett{\usebibmacro{date+extrayear}}}

\DefineBibliographyStrings{ngerman}{% Anpassungen für deutsche Sprache
	nodate = {{o.J.}},
	urlseen = {{Abruf:}},
	ibidem = {{ebenda}},
	andothers = {{et\addabbrvspace al\adddot}}
}
\DefineBibliographyStrings{english}{% Anpassungen für englische Sprache
    nodate = {{w.y.}},
    urlseen = {{retrieval:}}
}

% keine Anführungszeichen beim Titel im Literaturverzeichnis
\DeclareFieldFormat[article,book,inbook,inproceedings,manual,misc,phdthesis,thesis,online,report]{title}{#1\isdot}

\newcommand{\literaturverzeichnis}{%
% nur Literaturverzeichnis
% (als eigenes Kapitel)
\phantomsection
\addcontentsline{toc}{chapter}{\refname}
\spezialkopfzeile{\refname}
\defbibheading{lit}{\chapter*{\refname}}
\label{chapter:quellen}
\printbibliography[heading=lit,notkeyword=ausblenden]
}


% Verzeichnisse/Tabellen
\usepackage{tocloft}
\usepackage{multirow}
\usepackage{amsmath}
\usepackage{amssymb}
\usepackage{booktabs}

% Hyperlinks
\usepackage[hypertexnames=false]{hyperref}

% Bessere Umbrüche in \url{…}
\usepackage{xurl}
\urlstyle{tt}

% Intelligente Referenzen (nach hyperref laden!)
\usepackage[capitalise,nameinlink,noabbrev]{cleveref}
% "Listing" -> "Quellcode" (auch für Verzeichnis und Verweise)
\renewcommand{\lstlistingname}{Quellcode}
\renewcommand{\lstlistlistingname}{Quellcodeverzeichnis}
\providecommand*{\lstlistingautorefname}{Quellcode}
\crefname{lstlisting}{Quellcode}{Quellcode}
\Crefname{lstlisting}{Quellcode}{Quellcode}

% Anhangszähler/Verzeichnis (wie in deiner Vorlage)
\newcounter{anhcnt}\setcounter{anhcnt}{0}
\newlistof{anhang}{app}{}
\newcommand{\anhang}[1]{%
  \refstepcounter{anhcnt}\setcounter{anhteilcnt}{0}
  \section*{\appendixname\ \theanhcnt: #1}
  \addcontentsline{app}{section}{\protect\numberline{\appendixname\ \theanhcnt}#1}\par
}
\newcounter{anhteilcnt}\setcounter{anhteilcnt}{0}
\newcommand{\anhangteil}[1]{%
  \refstepcounter{anhteilcnt}
  \subsection*{\appendixname\ \arabic{anhcnt}/\arabic{anhteilcnt}: #1}
  \addcontentsline{app}{subsection}{\protect\numberline{\appendixname\ \theanhcnt/\arabic{anhteilcnt}}#1}\par
}
\renewcommand{\theanhteilcnt}{\appendixname\ \theanhcnt/\arabic{anhteilcnt}}

% tocloft-Einrückungen für Anhangverzeichnis
\makeatletter
\newcommand{\abstaendeanhangverzeichnis}{%
  \renewcommand*{\l@section}{\@dottedtocline{1}{0em}{5.5em}}
  \renewcommand*{\l@subsection}{\@dottedtocline{2}{2.3em}{6.5em}}
}
% Einträge LOF/LOT und Quellcodeverzeichnis (LOL) angeglichen
\renewcommand*{\l@figure}{\@dottedtocline{1}{0em}{2.3em}}
\renewcommand*{\l@table}{\@dottedtocline{1}{0em}{2.3em}}
\renewcommand*{\l@lstlisting}{\@dottedtocline{1}{0em}{2.3em}}
\makeatother

% Fortlaufende Zähler über Kapitel hinweg
\usepackage{chngcntr}
\counterwithout{figure}{chapter}
\counterwithout{table}{chapter}
\counterwithout{footnote}{chapter}

% Kopfzeilen (KOMA-Script)
\usepackage[automark]{scrlayer-scrpage}
%% Definitionen für Kopf- und Fußzeile auf normalen Seiten
\defpagestyle{kopfzeile}
{% Kopfdefinition
  (\textwidth,0pt)    % Länge der oberen Linie,Dicke der oberen Linie       
  {} % Definition für linke Seiten im doppelseitigen Layout
  {} % Definition für rechte Seiten im doppelseitigen Layout      
  {  % Definition für Seiten im einseitigen Layout
	\makebox[0pt][l]{\rightmark}% 
	\makebox[\linewidth]{}% 
  }        
  (\textwidth, 0.4pt) % Untere Linienlänge, Untere Liniendicke
}
{% Fußdefinition
  (\textwidth,0pt)    % Obere Linienlänge, Obere Liniendicke
  {} % Definition für linke Seiten im doppelseitigen Layout
  {} % Definition für rechte Seiten im doppelseitigen Layout
  {  % Definition für Seiten im einseitigen Layout
    \makebox[\linewidth]{}%
    \makebox[0pt][r]{\pagemark}%
  }
  (\textwidth, 0pt)   % Länge der unteren Linie,Dicke der unteren Linie
}

%% Definitionen für Kopf- und Fußzeile auf ersten Seiten eines Kapitels
\defpagestyle{kapitelkopfzeile}
{% Kopfdefinition
  (\textwidth,0pt)    % Länge der oberen Linie,Dicke der oberen Linie       
  {} % Definition für linke Seiten im doppelseitigen Layout
  {} % Definition für rechte Seiten im doppelseitigen Layout      
  {}  % Definition für Seiten im einseitigen Layout
  (\textwidth, 0pt) % Untere Linienlänge, Untere Liniendicke
}
{% Fußdefinition
  (\textwidth,0pt)    % Obere Linienlänge, Obere Liniendicke
  {} % Definition für linke Seiten im doppelseitigen Layout
  {} % Definition für rechte Seiten im doppelseitigen Layout
  {  % Definition für Seiten im einseitigen Layout
    \makebox[\linewidth]{}%
    \makebox[0pt][r]{\pagemark}%
  }
  (\textwidth, 0pt)   % Länge der unteren Linie,Dicke der unteren Linie
}

%% Definitionen für Kopf- und Fußzeile im Anhang und bei Quellenverzeichnisse
\newcommand{\spezialkopfzeileBezeichnung}{}
\defpagestyle{spezialkopfzeile}
{% Kopfdefinition
  (\textwidth,0pt)    % Länge der oberen Linie,Dicke der oberen Linie       
  {} % Definition für linke Seiten im doppelseitigen Layout
  {} % Definition für rechte Seiten im doppelseitigen Layout      
  {  % Definition für Seiten im einseitigen Layout
	\makebox[0pt][l]{\spezialkopfzeileBezeichnung}% 
	\makebox[\linewidth]{}% 
  }        
  (\textwidth, 0.4pt) % Untere Linienlänge, Untere Liniendicke
}
{% Fußdefinition
  (\textwidth,0pt)    % Obere Linienlänge, Obere Liniendicke
  {} % Definition für linke Seiten im doppelseitigen Layout
  {} % Definition für rechte Seiten im doppelseitigen Layout
  {  % Definition für Seiten im einseitigen Layout
    \makebox[\linewidth]{}%
    \makebox[0pt][r]{\pagemark}%
  }
  (\textwidth, 0pt)   % Länge der unteren Linie,Dicke der unteren Linie
}
            
\newcommand\spezialkopfzeile[1]{%
  \renewcommand\spezialkopfzeileBezeichnung{#1}
  \pagestyle{spezialkopfzeile}
}
                
% Standard-Pagestyle auswählen
\pagestyle{kopfzeile}

% keine Kopfzeile anzeigen auf Seiten, auf denen ein 
% Kapitel beginnt oder das Inhalts-/Abbildungs-/Tabellenverzeichnis steht 
\renewcommand{\chapterpagestyle}{kapitelkopfzeile}
\tocloftpagestyle{kapitelkopfzeile}



% Euro-Zeichen
\usepackage{textcomp}
\usepackage{eurosym}
\renewcommand{\texteuro}{\euro}  % ACHTUNG: nach hyperref laden!

% Kompatibilität KOMA-Script
\usepackage{scrhack}

% Abstände bei Kapitelüberschriften (inkl. Verzeichnisse)
\renewcommand*\chapterheadstartvskip{\vspace*{-\topskip}}
\newcommand{\myBeforeTitleSkip}{1mm}
\newcommand{\myAfterTitleSkip}{10mm}
\setlength\cftbeforetoctitleskip{\myBeforeTitleSkip}
\setlength\cftbeforeloftitleskip{\myBeforeTitleSkip}
\setlength\cftbeforelottitleskip{\myBeforeTitleSkip}
\setlength\cftaftertoctitleskip{\myAfterTitleSkip}
\setlength\cftafterloftitleskip{\myAfterTitleSkip}
\setlength\cftafterlottitleskip{\myAfterTitleSkip}

% Anhang beginnen
\newcommand{\startAnhang}{%
  \chapter*{\appendixname}
  \addcontentsline{toc}{chapter}{\appendixname}
  \section*{\anhangVzBezeichnung}
  \vspace{-8em}
  % vor \listofanhang müssen Einrückungen angepasst werden
  \abstaendeanhangverzeichnis
  \spezialkopfzeile{\DEoEN{Anhang}{Appendix}}
}

% Abkürzungsverzeichnis beginnen
\newcommand{\startAbkVerzeichnis}{%
  \chapter*{\abkVzBezeichnung}
  \addcontentsline{toc}{chapter}{\abkVzBezeichnung}
}

% Zeilenabstand in Tabellen schnell ändern
\newcommand{\ra}[1]{\renewcommand{\arraystretch}{#1}}

% Sprachspezifische Überschriften
\DEoEN{%
  \newcommand{\abkVzBezeichnung}{Abkürzungsverzeichnis}
  \newcommand{\anhangVzBezeichnung}{Anhangverzeichnis}
  \renewcaptionname{ngerman}{\refname}{Literaturverzeichnis}
  \renewcaptionname{ngerman}{\figurename}{Abb.}
  \renewcaptionname{ngerman}{\tablename}{Tab.}
}{%
  \newcommand{\abkVzBezeichnung}{Abbreviations}
  \newcommand{\anhangVzBezeichnung}{Appendix directory}
  \renewcaptionname{english}{\contentsname}{Table of Contents}
  \renewcaptionname{english}{\figurename}{Fig.}
  \renewcaptionname{english}{\tablename}{Tab.}
}

%%% Ende der Präambel %%%


%%% Name der eigenen Literatur-Datenbank (ggf. anpassen) %%%
\bibliography{includes/literatur-datenbank.bib}

\begin{document}
%%% Deckblatt gemäß DHBW-Vorgaben einbinden (keine Anpassung nötig) %%% 
    \begin{titlepage}
        \centering % Zentriert den gesamten Inhalt auf der Seite
       
        % Logos
        \includegraphics[height=60pt]{TruLogo_Print_RGB.jpg}\hfill\includegraphics[height=60pt]{DHBW_Logo.jpg}
       
        \vspace{1.5cm} % Abstand nach Logos
       
        % Titel
        {\LARGE\textsc{Erweiterung eines KI-gestützten Assistenzsystems zur Optimierung von
Laserschneidparametern für Edelstahlbleche}}\\[1.5cm]
       
        {\Large\textbf{Projektarbeit T2000}}\\[1cm]
       
        % Studiengang und Hochschule
        {\large
            des Studienganges Elektrotechnik\\
            Fachrichtung Automation\\
            an der Dualen Hochschule Baden-Württemberg\\
            Standort Stuttgart
        }\\[1.5cm]
       
        % Name des Autors
        {\large\textbf{Tudor Lupsa}}\\[0.5cm]
       
        % Abgabedatum
        {\large 08.09.2025}\\[2cm]
       
        % Weitere Informationen
        \begin{tabular}{ll}
            \textbf{Bearbeitungszeitraum} & 02.06.25 - 08.09.25\\
            \textbf{Matrikelnummer, Kurs} & 1491114, TEL23GR3\\
            \textbf{Dualer Partner} & TRUMPF SE+Co.KG, Ditzingen\\
            \textbf{Betreuer des Dualen Partners} & Manuel Geiger, M.Sc\\
    %       \textbf{Betreuer der Dualen Hochschule} & Name\\
        \end{tabular}
       
    \end{titlepage}


\pagenumbering{Roman}

%\begin{abstract}
\thispagestyle{kapitelkopfzeile}


\end{abstract}




%%% Inhalts-, Abbildungs-, Tabellenverzeichnisse %%%
% werden einzeilig gesetzt, um Platz zu sparen 
\begin{spacing}{1}
\tableofcontents % Inhaltsverzeichnis ausgeben
\clearpage
\startAbkVerzeichnis

Ein Abkürzungsverzeichnis ist optional. Das Paket \verb|acronym| kann weit mehr, als hier gezeigt.\footnote{siehe \url{http://ctan.org/pkg/acronym}}
Beachten Sie allerdings, dass Sie die Einträge selbst in sortierter Reihenfolge angeben müssen.

\begin{acronym}[DHBW] 
% Argument definiert die Breite der ersten Spalte anhand des längsten vorkommenden Eintrags
\acro{CRM}{Customer Relationship Management}
\acro{DHBW}{Duale Hochschule Baden-Württemberg}
\acro{IEEE}{Institute of Electrical and Electronics Engineers}
\acro{ITIL}{IT Infrastructure Library}
\acro{RoI}{Return-On-Invest}
\acro{UCS}{Universal Character Set}
\acro{UTF-8}{8-Bit UCS Transformation Format} 
\end{acronym}

\vspace{2em}
{\footnotesize
\textbf{Ergänzende Bemerkung:}
Eine im Text verwendete Abkürzung sollte bei ihrer ersten Verwendung erklärt werden. Falls Sie sich nicht selbst darum kümmern möchten, kann das das Paket \verb|acronym| übernehmen und auch automatisch Links zum Abkürzungsverzeichnis hinzufügen. Dazu ist an allen Stellen, an denen die Abkürzung vorkommt, \verb|\ac{ITIL}| zu schreiben. 

Das Ergebnis sieht wie folgt aus: 
\begin{itemize}
\item erstmalige Verwendung von \verb|\ac{ITIL}| ergibt: \ac{ITIL},
\item weitere Verwendung von \verb|\ac{ITIL}| ergibt: \ac{ITIL}
\end{itemize}
Wo benötigt, kann man mit dem Befehl \verb|\acl{ITIL}| wieder die Langfassung ausgeben lassen: \acl{ITIL}.

Falls man die Abkürzungen durchgängig so handhabt, kann man durch Paket-Optionen (in \verb|_dhbw_praeambel.tex|)
erreichen, dass im Abkürzungsverzeichnis nur die tatsächlich verwendeten Quellen aufgeführt werden (Option: \verb|printonlyused|) und zu jedem Eintrag die Seite der ersten Verwendung angegeben wird (Option: \verb|withpage|).

Durch die aktivierte Paket-Option \verb|nohyperlinks| wird verhindert, dass die Einträge im Abkürzungsverzeichnis mit Links zu der Stelle hinterlegt werden, wo der Begriff zum ersten Mal verwendet wird.} % Abkürzungsverzeichnis einbinden

\clearpage
\thispagestyle{kapitelkopfzeile}
\listoffigures
\phantomsection
\addcontentsline{toc}{chapter}{\listfigurename} % Abb.verz. ins Inh.verz. aufnehmen

\clearpage
\listoftables
\phantomsection
\addcontentsline{toc}{chapter}{\listtablename} % Tab.verz. ins Inh.verz. aufnehmen
\end{spacing}

%%% Umstellung der Seiten-Nummerierung auf 1, 2, 3 ... %%%
\cleardoublepage
\pagenumbering{arabic}

%%% Ihr eigentlicher Inhalt %%%
% Empfehlung: strukturieren Sie Ihren Text in einzelnen Dateien 
% und binden Sie diese hier mit \input{includes/dateiname.tex} ein
\chapter{Einführung}
\section{Zielsetzung}
Das Ziel dieser Projektarbeit ist es, ein bestehendes KI-gestütztes Assistenzsystem zu erweitern, das aktuell die Qualität beim Laserschneiden von Baustahlblechen vorhersagt und optimiert. Konkret soll die Leistungsfähigkeit dieses Modells auf Edelstahlbleche übertragen werden. Das aktuelle KI-Modell weist in Bezug auf Edelstahl Defizite auf, da es bisher nur mit Datensätzen von Baustahl trainiert wurde und die spezifischen Eigenschaften von Edelstahl unzureichend berücksichtigt.
Um diese Lücke zu schließen, sollen neue, speziell auf Edelstahl zugeschnittene Trainingsdaten generiert und in das bestehende Modell integriert werden. Diese Daten erfassen insbesondere typische Eigenschaften wie Schneidgratbildung und Oberflächenrauheit. Zusätzlich werden bestehende optische Messmethoden geprüft und gegebenenfalls angepasst, um ihre Eignung für Edelstahl sicherzustellen. Das Ziel ist ein robustes und zuverlässiges KI-Modell, das die Qualität von Edelstahlschnitten ebenso präzise vorhersagen kann wie bereits für Baustahl.

\section{Vorgehensweise}
Zunächst wird ein systematischer Testplan erstellt, um wichtige Schneidparameter, insbesondere Laserleistung, Schnittgeschwindigkeit, Gasdruck und Fokuslage, für verschiedene Edelstahldicken und -güten zu untersuchen. Ziel dieses Testplans ist es, Parameterbereiche zu identifizieren, die zu sogenannten Schnittabrissen führen. Ein Schnittabriss entsteht, wenn durch ungünstige Schneidparameter das Werkstück nicht vollständig getrennt wird, da die Schnittkante verschweißt.

Nachdem kritische Parameterbereiche, welche zu schlechten Schneidergebnissen führen, ermittelt wurden, werden detaillierte Versuchspläne („Experimentalplans“) entwickelt. Diese umfassen systematische Schneidversuche für Edelstahlbleche mit Dicken bis zu 20 mm. Hierbei werden bewusst sowohl qualitativ hochwertige als auch minderwertige Schneidergebnisse erzeugt. Dies dient dazu, eine umfassende Datengrundlage für die Weiterentwicklung des KI-Modells bereitzustellen.

Die generierten Schneidproben werden anschließend in einer Messzelle vermessen, um die resultierenden Schnittkanten in die KI-Datenbank aufzunehmen. Da die Messmethoden ursprünglich für Baustahl entwickelt wurden, müssen sie für Edelstahl angepasst werden. Dies betrifft insbesondere die Kalibrierung und Einrichtung des Handscanner-Setups, welches Bilder der Schnittkanten für die Qualitätsschätzung aufnimmt. Ebenso muss die Vektorberechnung des eingesetzten 3D-Punktwolkenscanners optimiert werden. Da Edelstahlschnittkanten typischerweise ausgeprägtere Schneidgrate aufweisen als die von Baustahlblechen, muss die Umrechnung der 3D-Punktwolke in einen 2D-Vektor entsprechend angepasst werden, um die tatsächlichen Merkmale der Schnittkanten präzise abzubilden.
Neben der quantitativen Messung des Schneidgrats wird auch die Oberflächenrauheit qualitativ bewertet. Hierzu werden bestehende Algorithmen zur Bildverarbeitung und Analyse geprüft und entsprechend den spezifischen Anforderungen von Edelstahl optimiert.
Die aufbereiteten Messdaten fließen anschließend in die Erweiterung und das Training des bestehenden KI-Modells ein. Ziel ist, dass dieses Modell anschließend die Qualität der Laserschneidkanten bei Edelstahlblechen zuverlässig vorhersagen kann. Nach erfolgreicher Implementierung erfolgen Validierungstests sowie weitere gezielte Optimierungen, um die Vorhersagequalität kontinuierlich zu verbessern und sicherzustellen, dass das gewählte Parameterset bereits vor dem Schneidprozess zuverlässig bewertet werden kann.
\chapter{Grundlagen}

\chapter{Stand der Technik}
\section{KI-gestütesztes Laserschneiden - Cutting Assistent}
\label{sec:cutting-assistent}

\section{3D-Messsystem - Keyence}
\label{sec:3d-messsystem-keyence}
\blinddocument
%%% Ende des eigentlichen Inhalts %%%


%%% Quellenverzeichnisse (keine Anpassung nötig) %%%
\clearpage
\literaturverzeichnis
%%% Ende Quellenverzeichnisse %%%


%%% Erklärung (keine Anpassungen nötig) %%%
% steht ganz am Ende des Dokuments
\cleardoublepage
\clearpage

\thispagestyle{empty}
\DEoEN{%
{\LARGE\textsf{\textbf{Erklärung zur Verwendung generativer KI-Systeme}}\bigskip}

Bei der Erstellung der eingereichten Arbeit habe ich auf künstlicher Intelligenz (KI) basierte Systeme benutzt: 

\begin{enumerate}
\item[$\boxtimes$] ja
\item[$\square$] nein\footnote{%
Die Erklärung ist in jedem Fall zu unterzeichnen, auch wenn Sie keine KI-Systeme genutzt haben und Ihr Kreuz bei \enquote{nein} gesetzt haben.}
\end{enumerate}

% Falls nein: Rest des Blocks entfernen, Unterschrift aber beibehalten

Falls ja: Die nachfolgend aufgeführten auf künstlicher Intelligenz (KI) basierten Systeme habe ich bei der Erstellung der eingereichten Arbeit benutzt:

\begin{enumerate}
\item
\item
\item \ldots
\end{enumerate}

Ich erkläre, dass ich

\begin{itemize}
  \item mich aktiv über die Leistungsfähigkeit und Beschränkungen der oben genannten 
KI-Systeme informiert habe,\footnote{U.a. gilt es hierbei zu beachten, dass an KI weitergegebene Inhalte ggf. als Trainingsdaten genutzt und wiederverwendet werden. Dies ist insb. für betriebliche Aspekte als kritisch einzustufen.}
  \item die aus den oben angegebenen KI-Systemen direkt oder sinngemäß übernommenen Passagen gekennzeichnet habe,
%
% In der Fußnote Ihrer Arbeit geben Sie die KI als Quelle an, z.B.: 
% Erzeugt durch Microsoft Copilot am dd.mm.yyyy. 
% Oder: Entnommen aus einem Dialog mit Perplexity vom dd.mm.yyyy. 
% Oder: Absatz 2.3 wurde durch ChatGPT sprachlich geglättet.
%
  \item überprüft habe, dass die mithilfe der oben genannten KI-Systeme generierten und von mir übernommenen Inhalte faktisch richtig sind,
  \item mir bewusst bin, dass ich als Autorin bzw. Autor dieser Arbeit die Verantwortung für die in ihr gemachten Angaben und Aussagen trage.
\end{itemize}

Die oben genannten KI-Systeme habe ich wie im Folgenden dargestellt eingesetzt: 


\begin{tabular}{|p{4cm}|p{3cm}|p{7cm}|}
    \hline
    \textbf{Arbeitsschritt in der wissenschaftlichen Arbeit} &
%
% Beispiele hierfür sind u.a. die folgenden Arbeitsschritte: 
% Generierung von Ideen, Konzeption der Arbeit, Literatursuche, Literaturanalyse, 
% Literaturverwaltung, Auswahl von Methoden, Datensammlung, Datenanalyse, 
% Generierung von Programmcodes
%
% Wenn Sie unsicher sind, ob Sie ein verwendetes KI-System angeben müssen, 
% wenden Sie sich an Ihre:n Betreuer:in.
%
    \textbf{Eingesetzte(s) KI-System(e)} & \textbf{Beschreibung der Verwendungsweise} \\
    \hline
    & & \\ % Platzhalter für Ihre Einträge, fügen Sie hier die Informationen ein
    \hline
    & & \\ % (Die Tabelle ist im Bedarfsfall zu erweitern und auf den Folgeseiten fortzusetzen)
    \hline
    & & \\
    \hline
    & & \\
    \hline
  \end{tabular}
} % Ende deutscher Teil
{% Beginn englische Erklaerung
{\LARGE\textsf{\textbf{Declaration on the Use of Generative AI Systems}}\bigskip}

In preparing the submitted work, I have used artificial intelligence (AI)-based systems:

\begin{enumerate}
\item[$\boxtimes$] yes
\item[$\square$] no\footnote{%
The declaration must be signed in any case, even if you have not used any AI systems and have checked the box \enquote{no}.}
\end{enumerate}

% Falls nein: Rest des Blocks entfernen, Unterschrift aber beibehalten

If yes: I have used the following AI-based systems:


\begin{enumerate}
\item
\item
\item \ldots
\end{enumerate}

I hereby declare that I

\begin{itemize}
  \item actively informed myself about the capabilities and limitations of the above-mentioned AI systems,\footnote{In particular, it should be noted that content passed on to AI may be used as training data and reused. This is to be considered critical, especially for operational aspects.}
  \item indicated passages directly or indirectly adopted from the above-mentioned AI systems,
%
% In the footnote of your work, cite the AI as the source, e.g.:
% Generated by Microsoft Copilot on mm.dd.yyyy. 
% Or: Taken from a dialogue with Perplexity on mm.dd.yyyy. 
% Or: Paragraph 2.3 was linguistically smoothed by ChatGPT.
%
  \item verified that the content generated and adopted by me using the above-mentioned AI systems is factually correct,
  \item am aware that as the author of this work, I bear responsibility for the statements and information provided in it.
\end{itemize}

I have utilized the above-mentioned AI systems as illustrated below: 

\begin{center}
\begin{tabular}{|p{4cm}|p{3cm}|p{7cm}|}
    \hline
    \textbf{Task in the scientific work} &
%
% Examples of such tasks include: idea generation, 
% conception of the work, literature search, literature analysis, 
% literature management, selection of methods, data collection, 
% data analysis, generation of program code
%
% If you are unsure whether you need to specify an AI system used, 
% consult your supervisor.
%
	 \textbf{AI System(s) Used} & \textbf{Description of Usage} \\
    \hline
    & & \\ % Placeholder for your entries, insert your information here
    \hline
    & & \\ % (Expand the table as needed and continue on subsequent pages)
    \hline
    & & \\
    \hline
    & & \\
    \hline
  \end{tabular}
\end{center}
}

\vspace{3cm}

\begin{center}
\begin{tabular}{ccc}
(\DEoEN{Ort, Datum}{place, date}) & \hspace{0.3\linewidth} & (\DEoEN{Unterschrift}{signature})
\end{tabular}
\end{center}
\clearpage

\thispagestyle{empty}

{\LARGE\textsf{\textbf{\DEoEN{Erklärung}{Declaration}}}\bigskip}

% \typMeinerArbeit und \themaMeinerArbeit werden in config.tex definiert
\DEoEN{%
Ich versichere hiermit, dass ich die vorliegende Arbeit mit dem Thema: \emph{\themaMeinerArbeit} selbstständig verfasst und keine anderen als die angegebenen Quellen und Hilfsmittel benutzt habe.
Ich versichere zudem, dass die eingereichte elektronische Fassung mit der gedruckten Fassung übereinstimmt.%
}{%
I hereby insure that I have personally authored this thesis with the topic: \emph{\themaMeinerArbeit} and have used no sources and aids other than those indicated. I also insure that the submitted electronic version corresponds to the printed version.%
}

\vspace{3cm}

\begin{center}
\begin{tabular}{ccc}
(\DEoEN{Ort, Datum}{place, date}) & \hspace{0.3\linewidth} & (\DEoEN{Unterschrift}{signature})
\end{tabular}
\end{center}


\end{document}
